\subsection{离散赋值}

定义 3. 设 K 是一个（未必交换的）域，v 是 K 上的一个赋值，\(\Gamma\) 是 v 的阶群。若存在由有序群 \(\Gamma\) 到 Z 上的一个（必然唯一的）同构，则称 v 为离散的。设 \(\gamma\) 是在此同构下对应于 1 的 \(\Gamma\) 的元素；K 中满足 \(v(u) = \gamma\) 的每个元素 u 都称为 v 的一个单值化元。若离散赋值的阶群为 Z，则称它是赋范的。

例如，由主理想整环*或因子分解整环*的极端元 p 所定义的赋值 \(v_{p}\) 是一个赋范离散赋值，它以 p 为单值化元。特别地，若 k 是域，则 \(k[[T]]\) 是 \(k((T))\) 上一个离散赋值的环，该赋值以 T 为单值化元。* 设 S 是一维连通复解析空间，K 是 S 上的亚纯函数域，\(z_{0}\) 是 S 的一点；\(f \in K\) 中在 \(z_{0}\) 处全纯的函数组成的集合是某个离散赋值 v 的环；函数 \(f \in K\) 是 v 的单值化元，当且仅当它在 \(z_{0}\) 处全纯且取值为零，并且存在 \(z_{0}\) 在 S 中的一个邻域 V，使 f 在 V 上的限制是 V 到 C 中原点的某邻域上的一个同胚。正是这个例子以及其他类似的例子成为“单值化元”一词的来源。*

命题 8. 设 K 是一个（未必交换的）域，v 是 K 上的一个离散赋值，A 是 v 的环，u 是 v 的一个单值化元。A 的非零理想都是双侧理想，且形如 Au \(^{n}\)（n ≥ 0）。

可以假定 \(v\) 是赋范的，于是 \(v(u) = 1\) 。对一切 \(x \in \mathbf{K}^*\) ，存在整数 \(n \in \mathbf{Z}\) 使得 \((v(x) = n = v(u^n))\)，从而可以写出

\[
x = z u ^ {n} = u ^ {n} z ^ {\prime},
\]

其中 \(z, z'\) 是环 A 的两个可逆元；由此即得命题。

命题 9. 设 A 是一个与其分式域不同的局部整环。下列条件等价：

\begin{enumerate}
\def\labelenumi{(\alph{enumi})}
\item
  A 是一个离散赋值的环。
\item
  A 是主理想整环。
\item
  理想 \(\mathfrak{m}(\mathbf{A})\) 是主理想，且 \(\bigcap_{n=1}^{\infty} \mathfrak{m}(\mathbf{A})^n = (0)\) 。
\item
  A 是 Noether 环，且 \(\mathfrak{m}(\mathbf{A})\) 是主理想。
\item
  A 是 Noether 赋值环。
\end{enumerate}

命题 8 表明 (a) 蕴含 (b)、(d) 与 (e)。若 A 是主理想整环，则由于 A 是局部环（《代数》第七章 §1 第 3 节定理 2），\(\mathfrak{m}(\mathbf{A}) = \mathbf{A}\mathbf{u}\)，且 A 的每个非零理想都形如 \(\mathbf{A}\mathbf{u}^n\)；因此 \(\bigcap_{n=1}^{\infty} \mathfrak{m}(A)^n = 0\) ；这表明 (b) 蕴含 (c)。另一方面，(d) 蕴含 (c)（第三章 §3 第 2 节命题 5 的推论）；由 §1 第 4 节命题 2，(c) 蕴含 (a)。于是条件 (a)、(b)、(c)、(d) 等价并且蕴含 (e)。最后假设 (e) 成立，我们来证明 (b) 成立；只需证明下面的引理：

引理 1. 设 A 是赋值环。每个有限生成的无扭 A-模都是自由模。A 的每个有限生成理想都是主理想。每个无扭 A-模都是平坦模。

设 E 是有限生成的无扭 A-模，\(x_{1}, \ldots, x_{n}\) 是 E 的一组个数最少的生成元；我们证明它们线性无关。若 \(\sum_{i=1}^{n} a_i x_i = 0\)（ \(a_i \in A\) ）是 \(x_i\) 之间的一个非平凡关系，则某个 \(a_i\)（例如 \(a_1\)）整除其余所有的元素，因为 \(A\) 的主理想集合按包含关系是全序的（§ 1 第 2 节定理 1）；又因该关系非平凡，故 \(a_1 \neq 0\)。由于 \(E\) 无扭，可以用 \(a_1\) 去除，这相当于假定 \(a_1 = 1\) 。但这样 \(x_1\) 就是 \(x_2, \ldots, x_n\) 的线性组合，与 \(n\) 的最小性相矛盾。因此 \(E\) 是自由的。

特别地，A 的每个有限生成理想 a 都是主理想，因为 a 的某个生成系中的所有元素都是其中一个的倍数。于是《代数》第一章 § 2 第 4 节命题 3 表明每个无扭 A-模都是平坦模。

